% 第八部：多數暴政（第陸部分，1395–1401）。出處：英文底稿；譯本附錄圖「決定上的多數與議題上的多數」。
\chapter{多數暴政}

\qaentry{第八部之一}{「暴政」一詞為何不能直接支持違憲審查}

\begin{qaquestion}
「多數暴政」是反對立法至上最常見的理由。作者如何處理「暴政」一詞？由多數所加諸，是否使暴政加重？
\end{qaquestion}

作者先指出，任何權利分歧的雙方都可以指對方的主張構成暴政；任何決定程序都可能作成錯誤的決定，法院也是\pg{1395--1396}。由多數所加諸，並不使暴政加重。他的看法是：多數決的一面反而減輕了暴政，因為該決定至少有一件事不是暴政式的，它沒有以暴政的方式把某些人排除在平等參與之外\pg{1396}。

\qaentry{第八部之二}{決定上的多數與議題上的多數}

\begin{qaquestion}
作者區分哪兩組多數與少數？表決落敗，何時才構成多數暴政？
\end{qaquestion}

作者區分「決定上的」（decisional）多數、少數與「議題上的」（topical）多數、少數\pg{1397}。該處的註說明，用 topical 一詞，是因為這些人的權利與利益正是該決定的議題；此一用語是寬鬆的，但足以讓人看出，並非每個在權利議題上投給落敗一方的人，都是權利受該決定不利影響的群體成員\pg{1397}。

作者以自己為例：我參與表決，我投的一方落敗，我因此屬於決定上的少數。要證明我遭受了暴政，還須證明兩件事：(1) 該決定確屬錯誤，對受影響者的權利帶有暴政的含意；(2) 我屬於權利受此錯誤決定不利影響的議題上的少數\pg{1397}。即使我屬於議題上的少數，情況也不必然改變，因為人們包括議題上的少數，未必擁有他們自認擁有的權利\pg{1398}。

\begin{qafigure}{決定上的多數、少數與議題上的多數、少數（1397--1401）。\label{fig:majorities}}
% 決定上與議題上的多數、少數（1397–1398、1401）。三格灰階，兩種少數重合的一格用強調色。
\begin{tikzpicture}[qa,cell/.style={nbox,text width=52mm,minimum height=15mm}]
  \node[hdr,align=center] at (3.15,1.25) {議題上的少數\\[-1pt]{\scriptsize\mdseries 權利與利益為該決定之議題者}};
  \node[hdr] at (8.85,1.25) {其他人};
  \node[hdr,anchor=east,align=right] at (0.15,0) {決定上的多數\\[-1pt]{\scriptsize\mdseries 表決勝方}};
  \node[hdr,anchor=east,align=right] at (0.15,-1.8) {決定上的少數\\[-1pt]{\scriptsize\mdseries 表決敗方}};
  \node[cell] at (3.15,0) {議題上的少數也有人投給勝方；墮胎權、平權措施的支持與反對者跨越性別與族群（1401）};
  \node[cell] at (8.85,0) {勝方若主要由議題上的多數組成，對應左下格的重合（1401）};
  \node[cell,kbox] at (3.15,-1.8) {兩種少數重合：決定若確屬錯誤，構成多數暴政；普遍如此即屬非核心情形（1397--1398）};
  \node[cell] at (8.85,-1.8) {只是表決落敗：自己的意見未被採行，本身不構成暴政（1397--1398）};
\end{tikzpicture}

\end{qafigure}

作者的小結：

\begin{quote}
To sum up, tyranny of the majority is possible. But the term should not be used simply to mark the speaker's disagreement with the outcome of a majority decision. The most fruitful way of characterizing tyranny of the majority is to say that it happens when topical minorities are aligned with decisional minorities.\pg{1398}
\end{quote}

依此界定，多數暴政是非核心情形的特徵：人們除自己的權利之外，甚少關心少數或個人的權利。作者強調這與第三項假設不相容，不應同時談論一個承諾權利的社會、又說多數暴政在其中是普遍的可能\pg{1398}。

\qaentry{第八部之三}{對 Dworkin「背景條件」說的回應}

\begin{qaquestion}
Dworkin 主張尊重權利是任何決定制度具正當性的背景條件。\\
作者如何回應？
\end{qaquestion}

Dworkin 在 Freedom's Law 主張，人民的權利若未受尊重，民主決定本身即屬暴政，因為尊重權利是任何政治決定制度具正當性的背景條件\pg{1398--1399}。註中說明 Dworkin 自己表示這不是支持違憲審查的論據。作者提出兩點\pg{1400}。

\begin{enumerate}
\item Dworkin 似乎認為，決定若涉及民主本身，就不必追問作成決定的程序。作者不同意：若關於多數決程序的決定以排除女性參與的程序作成，以平等為據的反對不會因爭點是多數決的正當性而失去資格；即使爭點就是程序，我們仍在乎程序價值。
\item 「民主程序只在彼此尊重權利的人之間具正當性」有兩種讀法：須持有並實踐對彼此權利的正確見解，或認真對待彼此的權利並真誠設法確認其內容。前者過強，沒有任何可想像的政治制度滿足它；作者接受後者，而依後者，本文所設想的社會已滿足這項前提。
\end{enumerate}

\qaentry{第八部之四}{任何人不得為自己案件之法官}

\begin{qaquestion}
有人援引 nemo iudex in sua causa，主張權利問題不應由人民自己終局決定。\\
作者如何回應？
\end{qaquestion}

作者說，幾乎任何可想像的決定規則，最後都有人就自己的案件作決定；除非上訴的鏈條無窮無盡，總有某人或某機關的決定是終局的，而該終局者總是就自己見解的可接受性作裁斷\pg{1400--1401}。這項論證若有意義，是指議題上的多數不應同時是決定上起決定作用的多數；兩者重合時確有正當的疑慮，若普遍如此則屬非核心情形。作者以墮胎與平權措施為例，指出這類重合在美國違憲審查常處理的案件中甚少發生：許多女性支持墮胎權，許多男性也支持，也有許多女性反對；平權措施的支持與反對者同樣跨越族群\pg{1401}。
